\chapter{\KGA: Knowability-Guided Adaptation in Practice}\label{ch:kga}

The preceding chapters establish when label-free adaptation is theoretically
knowable---and when it is not.  This chapter describes \KGA{}
(Knowability-Guided Adaptation), the concrete decision algorithm that
operationalises those results.  We present the design goals, the four families
of label-free evidence, the four certificate estimators with their exact
formulas, the decision policy and its false-adapt guarantee, the public
\texttt{kga} package API, a full deployment pseudocode, operational guidance,
the served HTTP interface, and an honest account of what \KGA{} does not
guarantee.

% -----------------------------------------------------------------------
\section{Design goals and mechanism-agnostic wrapper architecture}
\label{sec:kga-design}
% -----------------------------------------------------------------------

\KGA{} is a \emph{decision wrapper}, not a new adaptation mechanism.  Its
single responsibility is to answer, \emph{before} any irreversible update, the
three-way question: should this candidate adaptation $f_a$ replace the frozen
model $f_0$ on this unlabelled target batch?  The adapting mechanism
itself---entropy-minimising Tent \cite{wang2021tent}, sample-filtered EATA
\cite{niu2022eata}, sharpness-aware SAR \cite{niu2023sar}, or any other
algorithm---is a black box that supplies $f_a$; \KGA{} wraps it uniformly.

This design is dictated by three requirements.

\begin{enumerate}
  \item \textbf{Label-free.}  The gate must decide from evidence $Z =
    \phi(X_{1:m}, f_0, f_a)$ that is computable without any target label.  No
    target-label signals may enter at any step---neither for benefit estimation,
    nor for certificate construction.

  \item \textbf{Mechanism-agnostic.}  A laboratory may wish to evaluate several
    candidate adaptation strategies (Tent vs.\ EATA vs.\ SAR) under the same
    safety envelope.  By treating $f_a$ as an input, \KGA{} provides that
    shared envelope at the cost of one forward pass of $f_0$ and $f_a$ per
    decision.

  \item \textbf{Finite-sample valid.}  The false-adapt probability---the chance
    of deploying an $f_a$ that actually hurts---must be bounded at a
    user-specified level $\alpha$ even for small batches.  This is
    Theorem~\ref{thm:cert} from Chapter~\ref{ch:theory1}.
\end{enumerate}

The overall architecture (Figure~\ref{fig:kga-arch}) places \KGA{} between the
adaptation candidate and the live deployment loop.  Calibration data (labelled
validation scores) flows in once to build the calibration residuals needed for
the split-conformal path; thereafter each incoming unlabelled batch is processed
in three stages: evidence computation, certificate estimation, and the trichotomy
decision.

\begin{figure}[htbp]
  \centering
  \includegraphics[width=\textwidth]{fig_architecture.pdf}
  \caption{%
    \KGA{} system architecture.  Labelled calibration scores and an unlabelled
    test batch enter the pipeline.  The label-free evidence $Z$ is computed
    first; then a finite-sample certificate $\widehat\Delta \pm \varepsilon$ is
    formed at level $\alpha$; finally the trichotomy maps the certificate to one
    of three actions.  The adaptation mechanism (Tent, EATA, SAR, or any other)
    is a plug-in; \KGA{} wraps it without modification.
  }
  \label{fig:kga-arch}
\end{figure}

% -----------------------------------------------------------------------
\section{Label-free evidence computation}
\label{sec:kga-evidence}
% -----------------------------------------------------------------------

The evidence $Z$ is a vector of statistics computable from the calibration
scores $S_\mathrm{val}$ and the unlabelled test scores $S_\mathrm{test}$,
without any target labels.  The implementation lives in
\texttt{kga/evidence.py}: the dataclass \texttt{Evidence} and the function
\texttt{compute\_evidence(calib\_scores, test\_scores)}.  The signals fall into
four families, each tied to a theoretical result.

\paragraph{Drift (KS distance).}
For each detector column $j$, a two-sample Kolmogorov--Smirnov test
(\texttt{scipy.stats.ks\_2samp}) compares the calibration and test score
marginals.  The implementation stores:
\begin{itemize}
  \item \texttt{ks\_mean} --- mean per-detector KS statistic;
  \item \texttt{ks\_max} --- maximum per-detector KS statistic.
\end{itemize}
Large KS drift is the observable signature of a covariate shift to which
the gate can react (Theorem~\ref{thm:pos}).  When the marginal is
preserved but the informative signal is destroyed---the ``covert failure''
regime of the mixed-regime experiment---both statistics remain near zero and
the distribution is non-identifiable in the sense of
Theorem~\ref{thm:imp}.

\paragraph{Disagreement.}
The disagreement signal measures the heterogeneity of the test scores
across detectors:
\begin{equation}
  \texttt{disagree} \;=\; 1 - \overline{\rho},
  \label{eq:kga-disagree}
\end{equation}
where $\overline{\rho}$ is the mean pairwise Spearman rank correlation of the
test score columns (computed via \texttt{\_rank\_norm} then \texttt{np.corrcoef}
on the rank-normalised test matrix).  For a single detector, \texttt{disagree}
is defined to be zero.  A larger value indicates a bigger disagreement region
$D = \{x : f_0(x) \neq f_a(x)\}$, and thus greater potential for the
sign-of-difference to be identifiable
(Theorems~\ref{thm:disagree}--\ref{conj:gen}).

\paragraph{Entropy and confidence shift.}
The entropy signals measure the spread of the pooled score distributions.
\texttt{\_hist\_entropy} computes the Shannon entropy (in nats) of the
20-bin normalised histogram of a flattened score array:
\[
  H(S) \;=\; -\sum_k p_k \ln p_k.
\]
The confidence (decisiveness) signal \texttt{\_decisiveness} measures the mean
distance of the min-max-normalised scores from the indecision point $0.5$:
\[
  \texttt{conf}(S) \;=\; 2\,\E\!\left[\,\left|\frac{S - S_{\min}}{S_{\max} -
  S_{\min}} - 0.5\right|\right].
\]
The \texttt{Evidence} dataclass stores:
\begin{itemize}
  \item \texttt{calib\_entropy}, \texttt{test\_entropy} --- per-split entropies;
  \item \texttt{calib\_conf}, \texttt{test\_conf} --- per-split decisiveness;
  \item \texttt{entropy\_shift} = \texttt{test\_entropy} $-$
    \texttt{calib\_entropy};
  \item \texttt{conf\_shift} = \texttt{test\_conf} $-$ \texttt{calib\_conf}.
\end{itemize}
A large positive \texttt{conf\_shift} with a collapsing marginal is the
entropy-collapse warning sign familiar from online TTA monitoring.

\paragraph{Importance-weight effective sample size (ESS).}
The function \texttt{\_importance\_ess} models the calibration and test score
laws as univariate Gaussians and forms the Gaussian-ratio importance weight for
each test point $t_i$:
\[
  w_i \;=\; \frac{\mathcal{N}(t_i;\,\mu_c, \sigma_c^2)}
                  {\mathcal{N}(t_i;\,\mu_t, \sigma_t^2)},
\]
where $(\mu_c, \sigma_c)$ and $(\mu_t, \sigma_t)$ are the empirical moments of
the pooled calibration and test scores respectively.  The effective sample size
is then the standard self-normalised estimator,
\begin{equation}
  \mathrm{ESS} \;=\; \frac{(\sum_i w_i)^2}{\sum_i w_i^2},
  \label{eq:kga-ess}
\end{equation}
capped in $[1, n_\mathrm{test}]$.  The dataclass stores \texttt{ess} and the
scale-free ratio \texttt{ess\_frac} $= \mathrm{ESS}/n_\mathrm{test}$.  Low
\texttt{ess\_frac} indicates poor source-to-target support overlap, which
undermines importance-weighted certificates and is a quantitative diagnostic for
the support-overlap condition of Theorem~\ref{thm:pos}.

The fixed-order feature vector \texttt{Evidence.to\_vector()} concatenates
$[\texttt{ks\_mean},\ \texttt{ks\_max},\ \texttt{disagree},\
\texttt{entropy\_shift},\ \texttt{conf\_shift},\ \texttt{calib\_entropy},\
\texttt{test\_entropy},\ \texttt{calib\_conf},\ \texttt{test\_conf},\
\texttt{ess\_frac}]$ and can be passed directly to a downstream benefit
regressor as the feature vector $Z$.

\begin{remark}[Provenance]
  Every field of \texttt{Evidence} mirrors a quantity computed in the K-Bound
  experiment scripts: \texttt{kga/evidence.py} is a clean, importable API over
  the same mathematics, not a re-derivation.  The KS drift and disagreement
  fields match the $Z = \mathrm{dict}(\ldots)$ block in
  \texttt{src/scripts/kbound/knowability\_experiment.py} (lines~74--96); the
  entropy and importance-weight fields match the \texttt{evidence} helper in
  \texttt{mixed\_regime\_experiment.py}.
\end{remark}

% -----------------------------------------------------------------------
\section{Certificate estimators}
\label{sec:kga-certificates}
% -----------------------------------------------------------------------

A \emph{certificate} is a frozen dataclass \texttt{Certificate(delta\_hat,}
\texttt{epsilon, method, alpha, n)} implemented in \texttt{kga/certificate.py}.
It carries a point estimate $\widehat\Delta$ of the benefit
$\Delta = R(f_0) - R(f_a)$ and a confidence radius $\varepsilon$ such that, with
probability at least $1 - \alpha$,
\[
  \Delta \;\geq\; \widehat\Delta - \varepsilon \qquad \text{(certified lower bound)}.
\]
The two derived properties \texttt{lower} $= \widehat\Delta - \varepsilon$ and
\texttt{upper} $= \widehat\Delta + \varepsilon$ are the bounds consumed by the
decision policy.

The module provides four estimators, each linked to a theoretical result.

\subsection{Empirical Bernstein (\texttt{ebern}) --- the default}
\label{sec:kga-ebern}

The default estimator implements the Maurer--Pontil (2009) empirical-Bernstein
lower confidence bound \cite{maurer2009empirical}.  Given $n$ i.i.d.\ paired
benefits $X_i = \ell(f_0(x_i), y_i) - \ell(f_a(x_i), y_i)$ with empirical
mean $\widehat\mu$, unbiased sample variance $\widehat{V}$ (denominator $n-1$),
and range $R = b - a$, the certificate is:

\begin{equation}
  \varepsilon_{\mathrm{ebern}} \;=\;
  \sqrt{\frac{2\,\widehat{V}\,\ln(2/\alpha)}{n}}
  \;+\;
  \frac{7R\,\ln(2/\alpha)}{3(n-1)},
  \label{eq:ebern}
\end{equation}

so that $\Pr[\Delta \geq \widehat\mu - \varepsilon_{\mathrm{ebern}}] \geq 1 -
\alpha$.  The estimate is $\widehat\Delta = \widehat\mu$ and \texttt{epsilon} is
the sum of the two terms in~\eqref{eq:ebern}.  When $n = 1$ the variance is
undefined and \texttt{epsilon} is set to $+\infty$, forcing \textsc{abstain}.
When \texttt{benefit\_range} is not supplied, $R$ is estimated from the observed
range \texttt{max(X) - min(X)}.

This estimator is strictly tighter than Hoeffding whenever $\widehat{V} \ll R^2$
--- the usual regime in which most paired benefits agree in sign --- and is
deterministic, streamable, and finite-sample valid.  It is the batch Theorem~3
certificate of Theorem~\ref{thm:cert}.

\subsection{Hoeffding (\texttt{hoeffding}) --- distribution-free baseline}
\label{sec:kga-hoeffding}

The Hoeffding estimator replaces the variance term by the worst-case
variance-free bound:
\begin{equation}
  \varepsilon_{\mathrm{hoeff}} \;=\; R \sqrt{\frac{\ln(1/\alpha)}{2n}}.
  \label{eq:hoeff}
\end{equation}
This radius is always at least as large as the empirical-Bernstein one in the
low-variance regime, so it is provided as a conservative baseline and sanity
comparator.  The \texttt{method} field is set to \texttt{'hoeffding'}.

\subsection{Split-conformal (\texttt{conformal}) --- cross-task estimator}
\label{sec:kga-conformal}

The split-conformal estimator is the radius used in the main K-Bound experiments
(Theorem~\ref{thm:cert}, conformal variant; \cite{vovk2005algorithmic,
angelopoulos2023conformal}).  It takes a point estimate $\widehat\Delta$ (e.g.,
from a leave-one-out benefit regressor) and a vector of held-out calibration
residuals $r_i = |\widehat\Delta_i - \Delta_i|$ from $n_\mathrm{cal}$ tasks:
\begin{equation}
  \varepsilon_{\mathrm{conf}} \;=\; \mathrm{quantile}\!\left(r_{1:n},\, 1 -
  \alpha\right).
  \label{eq:conformal}
\end{equation}
By the exchangeability of the calibration split, $\Pr[|\widehat\Delta -
\Delta| \leq \varepsilon_{\mathrm{conf}}] \geq 1 - \alpha$ for a fresh query.
The function \texttt{conformal\_split(delta\_hat, calib\_residuals, alpha)} is
called in \texttt{knowability\_experiment.py} (\texttt{eps = np.quantile(resid,}
\texttt{1 - alpha)}) and in \texttt{mixed\_regime\_experiment.py}.

\subsection{Anytime e-value (\texttt{evalue}) --- Theorem~3b}
\label{sec:kga-evalue}

The anytime estimator implements two one-sided betting e-processes over an
ordered benefit stream \cite{ville1939etude,waudbysmith2024betting}:
\begin{align}
  E_t^+ &= \prod_{i \leq t} (1 + \lambda_i X_i), \qquad
  \text{tests } H_0^+: \Delta \leq 0,
  \label{eq:eplus}\\
  E_t^- &= \prod_{i \leq t} (1 + \nu_i(-X_i)), \qquad
  \text{tests } H_0^-: \Delta \geq 0,
  \label{eq:eminus}
\end{align}
with predictable truncated-aGRAPA bets
$\lambda_i = \mathrm{clip}(\widehat\mu_{i-1}/\widehat\sigma^2_{i-1},\;0,\;
\delta_{\mathrm{cap}}/|a|)$ (and symmetrically $\nu_i$) formed from
the running mean and second moment up to step $i-1$.  Each process is a
nonnegative supermartingale under its null; by \textbf{Ville's inequality}
\cite{ville1939etude}, $\Pr(\exists t: E_t \geq 1/\alpha) \leq \alpha$
simultaneously over all sample sizes with no multiplicity correction.

The returned certificate encodes the current anytime decision in the standard
$\widehat\Delta \pm \varepsilon$ form:
\begin{itemize}
  \item If $\log E_t^+ \geq \log(1/\alpha)$: emit $\varepsilon = 0$ with
    $\widehat\Delta > 0$, forcing \textsc{adapt};
  \item If $\log E_t^- \geq \log(1/\alpha)$: emit $\varepsilon = 0$ with
    $\widehat\Delta < 0$, forcing \textsc{freeze};
  \item Otherwise: emit $\varepsilon = |\bar{X}| + \epsilon_0$ (brackets zero
    around the running mean), forcing \textsc{abstain}.
\end{itemize}
This is Theorem~\ref{thm:evalue} (anytime certificate).  The implementation
mirrors \texttt{experiments/kbound/theory\_validation/val\_thm3\_evalue.py} and
\texttt{kga/kbound\_pkg/kbound/eprocess.py}.

% -----------------------------------------------------------------------
\section{Decision policy and the false-adapt guarantee}
\label{sec:kga-policy}
% -----------------------------------------------------------------------

The decision rule is a deterministic threshold on the certificate bounds,
implemented as the \texttt{decide(certificate, alpha)} function in
\texttt{kga/policy.py}.  The \texttt{Decision} enum has three members:
\texttt{ADAPT}, \texttt{FREEZE}, \texttt{ABSTAIN}.  As a \texttt{str} subclass,
each member compares equal to its string value and serialises cleanly to JSON.

The rule applies strict inequalities at the boundaries:

\begin{center}
\begin{tabular}{ll}
\toprule
Condition & Decision \\
\midrule
$\widehat\Delta - \varepsilon > 0$ & \textsc{adapt} \\
$\widehat\Delta + \varepsilon < 0$ & \textsc{freeze} \\
otherwise & \textsc{abstain} \\
\bottomrule
\end{tabular}
\end{center}

A certificate with $\widehat\Delta = \varepsilon$ (lower bound exactly zero)
\textsc{abstain}s rather than \textsc{adapt}s; this is the conservative choice
and matches the reference experiment scripts.  An infinite $\varepsilon$
(e.g., from a single-sample empirical-Bernstein certificate) always
\textsc{abstain}s.

\begin{theorem}[False-adapt guarantee; Theorem~\ref{thm:cert}]
  \label{thm:kga-false-adapt}
  The certificate is constructed so that $\Pr[\Delta \geq \widehat\Delta -
  \varepsilon] \geq 1 - \alpha$.  The \textsc{adapt} branch fires only when
  $\widehat\Delta - \varepsilon > 0$, which on the good event forces $\Delta >
  0$.  Therefore
  \[
    \Pr\!\left(\text{\textsc{adapt}} \text{ and } \Delta \leq 0\right) \;\leq\;
    \alpha.
  \]
  The symmetric statement bounds the false-freeze probability.  The anytime
  e-value variant (Theorem~\ref{thm:evalue}) upgrades this to hold
  simultaneously over all stream prefixes via Ville's inequality.
\end{theorem}

The proof follows identically from the non-identifiability regime (via
Theorem~\ref{thm:imp}): when evidence $Z$ cannot separate the two opposite-benefit
worlds, the certificate radius brackets zero and \KGA{} \textsc{abstain}s.  This
is mandatory, not merely conservative (see Conjecture~\ref{conj:gen} and the
discussion in Chapter~\ref{ch:theory1} on forced abstention).

% -----------------------------------------------------------------------
\section{The \texttt{kga} package API}
\label{sec:kga-api}
% -----------------------------------------------------------------------

The \texttt{kga} Python package (top-level directory \texttt{kga/}) exposes the
three-stage pipeline through the \texttt{KGA} facade class in \texttt{kga/kga.py}.

\paragraph{Constructor.}

\begin{verbatim}
KGA(alpha=0.1, method='ebern')
\end{verbatim}

\texttt{alpha} is the operating miscoverage level in $(0, 1)$; it bounds the
false-adapt probability via Theorem~\ref{thm:cert}.  \texttt{method} selects the
default batch certificate estimator when per-sample paired benefits are supplied:
\texttt{'ebern'} (empirical Bernstein; the default), \texttt{'hoeffding'}, or
\texttt{'evalue'} (anytime Theorem~3b).  A \texttt{KGA} instance carries no
learned state, so a single instance is reusable and thread-safe.

\paragraph{Methods.}

\begin{itemize}
  \item \textbf{\texttt{.evidence(calib, test, **kwargs) -> Evidence}} ---
    computes the label-free evidence $Z$ by delegating to
    \texttt{compute\_evidence}; stores the result in \texttt{last\_evidence}.

  \item \textbf{\texttt{.certify(...) -> Certificate}} --- builds $\widehat\Delta
    \pm \varepsilon$.  Three calling conventions are supported:
    \begin{enumerate}
      \item \texttt{certify(scores=benefits)} --- per-sample paired benefits,
        processed by the batch estimator selected at construction (or overridden
        via \texttt{method=});
      \item \texttt{certify(adapt\_risk=..., freeze\_risk=...,}
        \texttt{calib\_residuals=...)} --- two scalar risks plus held-out
        residuals; uses the split-conformal radius;
      \item \texttt{certify(delta\_hat=..., calib\_residuals=...)} --- explicit
        point estimate plus residuals for split-conformal.
    \end{enumerate}
    The result is stored in \texttt{last\_certificate}.

  \item \textbf{\texttt{.decide(certificate=None) -> Decision}} --- applies the
    trichotomy to a certificate (defaults to \texttt{last\_certificate}); stores
    the result in \texttt{last\_decision}.

  \item \textbf{\texttt{.explain() -> dict}} --- returns all intermediate
    quantities (\texttt{alpha}, \texttt{method}, \texttt{evidence},
    \texttt{certificate}, \texttt{decision}) as a JSON-serialisable dict.
\end{itemize}

\paragraph{Quickstart.}
The following usage snippet is taken directly from \texttt{kga/README.md}:

\begin{verbatim}
import numpy as np
from kga import KGA

rng = np.random.default_rng(0)
kga = KGA(alpha=0.1, method="ebern")

# (1) Label-free evidence Z.
calib = rng.normal(0.0, 1.0, size=(500, 3))
test  = rng.normal(0.0, 1.0, size=(500, 3))
z = kga.evidence(calib, test)
print(z.ks_mean, z.disagree, z.ess_frac)

# (2) Certificate from per-sample paired benefits.
benefits = rng.normal(0.3, 0.1, size=400)
cert = kga.certify(scores=benefits, benefit_range=2.0)
print(cert.delta_hat, cert.epsilon, cert.lower)

# (3) Trichotomy decision.
print(kga.decide(cert))     # -> Decision.ADAPT

# (4) Audit.
import json; print(json.dumps(kga.explain(), indent=2))
\end{verbatim}

\paragraph{Command-line interface.}
The \texttt{kga.cli} module (entry point \texttt{python -m kga}) exposes a
\texttt{decide} subcommand:

\begin{verbatim}
python -m kga decide --calib calib.npy --test test.npy --alpha 0.1
\end{verbatim}

This loads calibration and test scores from \texttt{.npy} files, computes the
evidence $Z$, constructs a split-conformal certificate with a label-free
residual proxy (\texttt{delta\_hat = 0}, radius from the calibration's own
dispersion), and prints the decision and certificate as JSON to stdout.  Because
labels are unavailable at the command line, it should be used for pipeline
integration testing rather than live deployment decisions; real benefit estimates
require the \texttt{KGA} API with paired benefits or held-out residuals.

% -----------------------------------------------------------------------
\section{Full deployment pseudocode}
\label{sec:kga-pseudocode}
% -----------------------------------------------------------------------

Algorithm~1 below gives the complete deployment loop, including the one-time
calibration phase and the per-batch inference phase.

\begin{center}
\fbox{\parbox{0.95\textwidth}{\small
\textbf{Algorithm 1: \KGA{} Deployment Loop.}\\[4pt]
\textbf{Inputs:} frozen model $f_0$; adaptation candidate $f_a$; calibration
split $(S_\mathrm{val}, y_\mathrm{val})$; unlabelled stream of target batches
$\mathcal{B}_1, \mathcal{B}_2, \ldots$; miscoverage level $\alpha$; certificate
method $m \in \{\texttt{ebern}, \texttt{hoeffding}, \texttt{evalue}\}$.\\[4pt]

\textbf{Calibration phase (once):}
\begin{enumerate}
  \item Compute paired benefits $X_i = \ell(f_0(x_i), y_i) - \ell(f_a(x_i), y_i)$
    on the calibration split.
  \item Fit a benefit estimator $\widehat\Delta(\cdot)$ (e.g., a leave-one-task-out
    gradient-boosted regressor) on $(Z_i, X_i)$ pairs.
  \item Record calibration residuals $r_i = |\widehat\Delta(Z_i) - \bar X_i|$
    for the split-conformal path, or retain $\{X_i\}$ for the batch path.
\end{enumerate}

\textbf{Per-batch inference phase (for each $\mathcal{B}_t$):}
\begin{enumerate}
  \setcounter{enumi}{3}
  \item \textbf{Evidence:} compute $Z_t = \phi(S_\mathrm{val}, S_t, f_0, f_a)$
    via \texttt{kga.evidence(calib, test)}.  Extract
    $[\texttt{ks\_mean},\, \texttt{ks\_max},\, \texttt{disagree},\,
    \texttt{entropy\_shift},\, \texttt{conf\_shift},\, \texttt{ess\_frac}]$.
  \item \textbf{Certify:} obtain $\widehat\Delta_t$ from the benefit estimator;
    call \texttt{kga.certify(scores=\ldots)} or
    \texttt{kga.certify(delta\_hat=}$\widehat\Delta_t$\texttt{,
    calib\_residuals=r)} to form $\varepsilon_t$ at level $\alpha$.
  \item \textbf{Decide (trichotomy):}
    \begin{itemize}
      \item \textbf{if} $\widehat\Delta_t - \varepsilon_t > 0$:
        $\Rightarrow$ \textbf{\textsc{adapt}} --- deploy $f_a$ on $\mathcal{B}_t$;
      \item \textbf{elif} $\widehat\Delta_t + \varepsilon_t < 0$:
        $\Rightarrow$ \textbf{\textsc{freeze}} --- keep $f_0$ on $\mathcal{B}_t$;
      \item \textbf{else}:
        $\Rightarrow$ \textbf{\textsc{abstain}} --- keep $f_0$, log for human
        review.
    \end{itemize}
  \item \textbf{Log:} write $(Z_t, \widehat\Delta_t, \varepsilon_t,
    \texttt{decision}_t, t)$ to the audit trail (\texttt{kga.explain()}).
\end{enumerate}
\textbf{Guarantee (Theorem~\ref{thm:cert}):} $\Pr(\textsc{adapt}\ \text{and}\ \Delta \leq 0)
\leq \alpha$; similarly for false-freeze.
}}
\end{center}

% -----------------------------------------------------------------------
\section{Operational guidance}
\label{sec:kga-ops}
% -----------------------------------------------------------------------

\paragraph{Choosing $\alpha$.}
The miscoverage level $\alpha$ directly bounds the false-adapt probability.  The
paper reports all results at $\alpha = 0.10$; production deployments in
safety-critical contexts should lower $\alpha$ (e.g., $0.05$), accepting a
higher abstention rate in exchange for stronger false-adapt control.  At very
small $\alpha$, the empirical-Bernstein radius grows as $\sqrt{\ln(2/\alpha)/n}$,
so the batch size should be increased proportionally to maintain coverage.

\paragraph{Batch size.}
The empirical-Bernstein and Hoeffding radii shrink as $O(1/\sqrt{n})$;
the split-conformal radius is independent of the batch size (it is a quantile
over the calibration tasks).  For the batch certificate to commit to
\textsc{adapt} or \textsc{freeze} rather than \textsc{abstain}, the benefit
$|\widehat\Delta|$ must exceed the radius.  In the 123-task anomaly-routing
benchmark, batch sizes of 200--500 samples provide useful separation at
$\alpha = 0.10$.

\paragraph{Calibration data.}
The calibration split must be exchangeable with deployment in the sense that the
per-task benefit distribution is representative (Theorem~\ref{thm:cert}
assumption~(i)).  In the cross-task anomaly-routing setting, a leave-one-task-out
procedure provides unbiased residuals.  In single-task streaming settings, a
held-out validation probe with the same label distribution as the deployment
stream is sufficient.

\paragraph{When \textsc{abstain} should trigger operational action.}
\textsc{abstain} means the evidence $Z$ is insufficient to certify the sign of
$\Delta$ at level $\alpha$; it does not mean that adapting is dangerous.  The
correct operational response is to (a) continue with the safer $f_0$ for the
current batch, (b) log the evidence vector and certificate for monitoring, and
(c) consider collecting more data or lowering the adaptation aggressiveness.
Repeated \textsc{abstain} outcomes on a stable stream may indicate that the
benefit estimator $\widehat\Delta$ is poorly calibrated for that domain, which
requires recalibration.

\paragraph{Complexity.}
\KGA{} is $O(n)$ per batch in the sample size $n$ and $O(k^2)$ in the number of
detectors $k$ (due to the pairwise rank-correlation in the disagreement
computation).  The implementation is pure \texttt{numpy}/\texttt{scipy},
deterministic (no random state), and torch-free, so it adds negligible overhead
to any deployment that already runs $f_0$ and $f_a$ forward passes.

% -----------------------------------------------------------------------
\section{The served HTTP interface}
\label{sec:kga-http}
% -----------------------------------------------------------------------

The production API serves \KGA{} decisions over HTTP through the routes defined
in \texttt{deploy/api/kga\_routes.py}, wired into the main FastAPI application.
Two endpoints are available.

\paragraph{\texttt{GET /kga/health} (no authentication).}
A liveness probe that returns the \texttt{kga} package version and component
status:

\begin{verbatim}
{"status": "ok", "component": "kga", "version": "0.1.0"}
\end{verbatim}

\paragraph{\texttt{POST /decide} (authentication required).}
Accepts a JSON body with three fields:
\begin{itemize}
  \item \texttt{calib\_scores} (\texttt{list[float]}, 2--200\,000 entries) ---
    calibration detector scores (label-free);
  \item \texttt{test\_scores} (\texttt{list[float]}, 2--200\,000 entries) ---
    unlabelled test detector scores;
  \item \texttt{alpha} (\texttt{float} in $(0,1)$, default $0.10$) ---
    miscoverage level.
\end{itemize}
All non-finite inputs are rejected with \texttt{HTTP 400}.  The response carries
the decision (\texttt{ADAPT}/\texttt{FREEZE}/\texttt{ABSTAIN}),
\texttt{delta\_hat}, \texttt{epsilon}, \texttt{method}, and the full
\texttt{Evidence} summary (\texttt{ks\_mean}, \texttt{ks\_max},
\texttt{disagree}, \texttt{entropy\_shift}, \texttt{conf\_shift},
\texttt{ess\_frac}, \texttt{n\_calib}, \texttt{n\_test}, \texttt{n\_detectors}).

Because target labels are unavailable at request time, the route reports a
conservative $\widehat\Delta = 0$ with a split-conformal radius derived from the
calibration scores' own dispersion (the same conservative proxy used in the
\texttt{kga.cli} entry point).  Callers with a real per-sample benefit estimate
should use the \texttt{KGA} Python API server-side rather than the HTTP route.

Request sizes are bounded at \texttt{MAX\_KGA\_SCORES = 200\,000} per array (the
same limit as \texttt{MAX\_*} in \texttt{deploy/api/main.py}).  The decision
math is entirely in the importable \texttt{kga} package; the route is a thin,
validated transport layer.

Figure~\ref{fig:kga-flow} shows the decision flow from an incoming HTTP request
to the returned decision.

\begin{figure}[htbp]
  \centering
  \includegraphics[width=\textwidth]{fig_decision_flow.png}
  \caption{%
    Decision flow for an incoming \texttt{POST /decide} request.  Pydantic
    validates and coerces the inputs; \texttt{compute\_evidence} produces $Z$;
    \texttt{conformal\_split} builds $\widehat\Delta \pm \varepsilon$; and
    \texttt{decide} applies the trichotomy.  Any \texttt{ValueError} from
    precondition checks (e.g., all-identical scores) is returned as \texttt{HTTP
    400}; uncaught exceptions return \texttt{HTTP 500}.
  }
  \label{fig:kga-flow}
\end{figure}

% -----------------------------------------------------------------------
\section{Threats, limitations, and what \KGA{} does not guarantee}
\label{sec:kga-threats}
% -----------------------------------------------------------------------

We state explicitly the situations in which \KGA{}'s guarantees do not apply,
to prevent misuse.

\paragraph{The false-adapt bound is not a guarantee that $f_a$ is good.}
Theorem~\ref{thm:cert} bounds $\Pr(\textsc{adapt}\ \text{and}\ \Delta \leq 0)
\leq \alpha$.  It says nothing about the magnitude of $\Delta$ when
$\Delta > 0$.  A deployment that \textsc{adapt}s on an $f_a$ with tiny
positive benefit has not violated the certificate; it has simply not gained much.
The certificate is about safety, not about optimality.

\paragraph{Non-identifiability is permanent, not a data problem.}
When two target worlds induce the same evidence distribution $Z$ but opposite
benefits (Theorem~\ref{thm:imp}), no amount of additional \emph{unlabelled}
data can resolve the ambiguity.  \KGA{} is forced to \textsc{abstain} in such
worlds by Corollary~\ref{sec:t2-cor-abstain}.  Gathering labelled target data
would resolve it, but that is outside the label-free scope.  The general
bracketing conjecture (Conjecture~\ref{conj:gen}) remains open.

\paragraph{The calibration exchangeability assumption.}
The split-conformal certificate (Section~\ref{sec:kga-conformal}) requires the
calibration tasks to be exchangeable with the deployment task.  If the
deployment stream is drawn from a radically different domain---e.g., calibrating
on anomaly-routing benchmarks and deploying on medical imaging---the certificate
radius $\varepsilon$ may be miscalibrated and the false-adapt bound may not hold.
This is the standard limitation of any conformal method.

\paragraph{Covert failures are not detectable.}
In the covert-failure regime of the mixed-regime experiment (the ``permuted
detector'' condition), the observable evidence $Z$ is nearly unchanged while the
true benefit reverses sign.  \KGA{} responds by widening the certificate radius
and \textsc{abstain}ing more frequently, which is correct given the theory, but
it means that under truly covert failures neither \textsc{adapt} nor
\textsc{freeze} can be certified.  Operational systems should treat an unusual
increase in \textsc{abstain} frequency as a sensor-quality warning.

\paragraph{The benefit estimator $\widehat\Delta$ must be calibrated.}
The false-adapt bound passes through the certificate; the certificate is valid
only insofar as $\varepsilon$ covers the estimation error $|\widehat\Delta -
\Delta|$.  If the benefit regressor is overfit to the calibration tasks or
trained on a distribution that does not include the deployment domain, the
effective coverage will be lower than $1 - \alpha$.  Regular recalibration
with fresh held-out data is therefore an operational requirement.

\paragraph{Label-free adaptation in the helpful-dominated regime.}
Where adaptation almost always helps---e.g., the clean 123-task suite where
$\text{mean}|\Delta| = 0.132$ on acted tasks---always-adapt is already a strong
baseline and the certificate's distinctive value is bounded by the frequency of
detectable harmful regimes.  We are explicit about this scope following the
honest-results principle articulated in the paper's scope note and in the
experimental discussion (Chapter~\ref{ch:exp1}).
Theorem~\ref{thm:imp} characterizes precisely the regime where this matters:
when the harmful-adaptation distribution has an observable fingerprint (large KS
drift, low ESS), \KGA{} can and does refuse it; when it does not
(covert failure), the Conjecture~\ref{conj:gen} bound suggests fundamentally
limited discriminability.